% \iffalse
%<*driver>
\ProvidesFile{seqenseigne-data.dtx}%
[2025/10/24 v1.0-dev seqenseigne – Module data (chargement & accès aux données)]
%</driver>
% \fi
%
% ============================================================================
% seqenseigne-data.dtx
% Module : DATA
% Rôle :
% - charger les CSV fusionnés (Objectifs.csv, Connaissances.csv)
% via datatool
% - filtrer par niveau / séquence courants
% - fournir des macros d'accès du type :
% \seqObjectifGetNom[niveau=N11,sequence=S01,objectif=03]
% \seqConnaissanceGetNom[...]
% - générer les tableaux d'objectifs:
% \seqTableauObjectifs
% ============================================================================
%
%<*package>
\NeedsTeXFormat{LaTeX2e}%
\ProvidesPackage{seqenseigne-data}%
[2025/10/24 v1.0-dev seqenseigne data module]

% On suppose que le core a déjà chargé les paquets génériques
% --- Bases de données en fichiers
\RequirePackage{datatool}%
\RequirePackage{xkeyval}%
\RequirePackage{xifthen}%
\RequirePackage{tabularray}%
%
% ============================================================================
% Variables internes
% ============================================================================
% -- Répertoire relatif des bases de données
\def\seq@var@datapath{../../}%
% -- Cache pour filtrer les données correctement par niveau et par séquence
\def\seq@var@cycle{C00}%
\def\seq@var@niveau{N00}%
\def\seq@var@sequence{S00}%
% -- Bases objectifs et connaissances filtrées par niveau et par séquence
% Bases par année du cycle (nommées anneeun/deux/trois, fixes)
\DTLnewdb{seq@objectifs@anneeun}%
\DTLnewdb{seq@objectifs@anneedeux}%
\DTLnewdb{seq@objectifs@anneetrois}%
\DTLnewdb{seq@objectifs@sequence}%
\DTLnewdb{seq@connaissances@anneeun}%
\DTLnewdb{seq@connaissances@anneedeux}%
\DTLnewdb{seq@connaissances@anneetrois}%
\DTLnewdb{seq@connaissances@sequence}%
% Bases par séquence : nommées seq@objectifs@sN (N = numéro décimal).
% Créées par boucle : ajuster seq@maxsequences pour supporter plus de séquences.
\newcounter{seq@tmp@i}%
\newcommand{\seq@maxsequences}{16}%
\setcounter{seq@tmp@i}{1}%
\@whilenum\value{seq@tmp@i}<\numexpr\seq@maxsequences+1\relax\do{%
	\DTLnewdb{seq@objectifs@s\arabic{seq@tmp@i}}%
	\stepcounter{seq@tmp@i}%
}%
%
% ----------------------------------------------------------------------------
% \begin{macro}{\seq@seqcode@tonumber}%
% Convertit un code séquence (S01, S02 … S16, S09, etc.) en son numéro décimal.
% Supprime le préfixe "S" et les zéros non significatifs en passant par \number.
% Résultat stocké dans la macro passée en #1.
% #1 : macro de sauvegarde (ex. \tmp@num)
% #2 : code séquence (ex. S01, S12)
% Usage interne uniquement (\seq@...).
\newcommand{\seq@seqcode@tonumber}[2]{%
	% On retire le "S" en développant le suffixe numérique brut (01, 02…)
	% puis \number supprime les zéros de tête : 01 → 1, 12 → 12.
	\edef\seq@tmp@stripped{\expandafter\@gobble\expandafter\@empty#2}%
	\edef#1{\number\seq@tmp@stripped}%
}%
%
% ============================================================================
% Commandes publiques
% ============================================================================
% \begin{macro}{\seqSetDataPath}%
% \IndexMacro{\seqSetDataPath}%
% Règle la valeur du chemin (relatif ou absolu) vers les bases de données
% Il doit se terminer par un /
% Le chemin par défaut est ../../../ pour compiler les livrets de séquence
% #1 : chemin
% \begin{macrocode}%
\newcommand{\seqSetDataPath}[1]{%
	\gdef\seq@var@datapath{#1}%
}%
% \end{macrocode}%
% \end{macro}%
% -------------------------------------------------------------------------------------------
% \begin{macro}{\seqLoadData}%
% \IndexMacro{\seqLoadData}%
% Les bases .dbtex sont créées avec l'outil externe datatooltk-gui 
% à partir des fichiers csv en prenant soin de régler "escape character" à "none".
% Elles contiennent, par cycle de 3 ans, les données concernant:
% - les themes: nombres et calculs, traitement des données, etc
% - les sequences: représentation d'un nombre, comaparaison de nombres, etc
% - les attendus de fin de cycle précédent répartis sur les séquences du cycle courant (uniquement pour le cycle 4), 
% - les connaissances de chaque séquence pour chaque niveau du cycle
% - les objectifs de chaque séquence pour chaque niveau du cycle
% \begin{macrocode}%
\newcommand{\seqLoadData}{%
	\input{\seq@var@datapath param_niveaux.dbtex}%
	\input{\seq@var@datapath C03_themes.dbtex}%
	\input{\seq@var@datapath C03_sequences.dbtex}%
	\input{\seq@var@datapath C03_connaissances.dbtex}%
	\input{\seq@var@datapath C03_objectifs.dbtex}%
	\input{\seq@var@datapath C04_themes.dbtex}%
	\input{\seq@var@datapath C04_sequences.dbtex}%
	\input{\seq@var@datapath C04_attendus_fin_C03.dbtex}%
	\input{\seq@var@datapath C04_connaissances.dbtex}%
	\input{\seq@var@datapath C04_objectifs.dbtex}%
}%
% \end{macrocode}%
% \end{macro}%
% -------------------------------------------------------------------------------------------
% \begin{macro}{\seqLoadDataSuivi}%
% \IndexMacro{\seqLoadDataSuivi}%
% #1 : code niveau
% \begin{macrocode}%
\newcommand{\seqLoadDataSuivi}[1]{%
	\DTLloaddb[keys={Numero}]{eleves}{liste_eleves.csv}%
	\DTLloaddb{periodes}{periodes.csv}%
	\DTLforeach*{periodes}{\Periode=Periode}{%
		\DTLloaddb{sequencesPeriode\Periode}{\Periode -sequences.csv}%
		\DTLloaddb[keys={Numero}]{bilanPeriode\Periode}{\Periode bilan.csv}%
	}%
	\input{param_niveaux.dbtex}%
	\seqDTLfetchsave{\tmp@codeCycle}{param_niveaux}{Code}{#1}{CodeCycle}%
	\ifthenelse{\equal{\tmp@codeCycle}{C03}}{%
		\input{C03_themes.dbtex}%
		\input{C03_sequences.dbtex}%
		\input{C03_objectifs.dbtex}%
		\input{C03_connaissances.dbtex}%
		\gdef\seq@var@cycle{C03}%
		\seq@bases@filter@cycle{C03}{N07}{N08}{N09}%
		\DTLforeach*{C03_sequences}{\tmp@codesequence=Code}{%
			\DTLloaddb[keys={Numero}]{Suivi\tmp@codesequence}{./\tmp@codesequence Suivi.csv}%
		}%
	}{%
		\ifthenelse{\equal{\tmp@codeCycle}{C04}}{%
			\input{C04_themes.dbtex}%
			\input{C04_sequences.dbtex}%
			\input{C04_objectifs.dbtex}%
			\input{C04_connaissances.dbtex}%
			\gdef\seq@var@cycle{C04}%
			\seq@bases@filter@cycle{C04}{N10}{N11}{N12}%
			\DTLforeach*{C04_sequences}{\tmp@codesequence=Code}{%
				\DTLloaddb[keys={Numero}]{Suivi\tmp@codesequence}{./\tmp@codesequence Suivi.csv}%
			}%
		}{% Code cycle inconnu : générer une erreur
		}%
	}%
}%
% \end{macrocode}%
% \end{macro}%
% -------------------------------------------------------------------------------------------
% \begin{macro}{\seqDTLfetchsave}%
% \IndexMacro{\seqDTLfetchsave}%
% Récupère et sauvegarde dans une macro la valeur d'un champ
% correspondant à la valeur d'un champ connu dans une base 
% #1 : nom de la macro dans laquelle on stocke le résultat
% #2 : base dans laquelle on recherche
% #3 : nom du champ de référence
% #4 : valeur du champ de référence
% #5 : nom du champ qu'on recherche
% \begin{macrocode}%
\newcommand{\seqDTLfetchsave}[5]{%
	\edtlgetrowforvalue{#2}{\dtlcolumnindex{#2}{#3}}{#4}%
	\dtlgetentryfromcurrentrow{\dtlcurrentvalue}{\dtlcolumnindex{#2}{#5}}%
	\let#1\dtlcurrentvalue
}%
% \end{macrocode}%
% \end{macro}%
% -------------------------------------------------------------------------------------------
% \begin{macro}{\seqNiveauGetNomCourt}%
% \IndexMacro{\seqNiveauGetNomCourt}%
% Renvoie le nom court du niveau spécifié
% #1 : code niveau
% \begin{macrocode}%
\newcommand{\seqNiveauGetNomCourt}[1]{%
	\DTLfetch{param_niveaux}{Code}{#1}{NomCourt}%
}%
% \end{macrocode}%
% \end{macro}%
% -------------------------------------------------------------------------------------------
% \begin{macro}{\seqNiveauGetNomLong}%
% \IndexMacro{\seqNiveauGetNomLong}%
% Renvoie le nom long du niveau spécifié
% #1 : code niveau
% \begin{macrocode}%
\newcommand{\seqNiveauGetNomLong}[1]{%
	\DTLfetch{param_niveaux}{Code}{#1}{NomLong}%
}%
% \end{macrocode}%
% \end{macro}%
% -------------------------------------------------------------------------------------------
% \begin{macro}{\seqNiveauGetCodeCycle}%
% \IndexMacro{\seqNiveauGetCodeCycle}%
% Renvoie le code du cycle du niveau spécifié
% #1 : code niveau
% \begin{macrocode}%
\newcommand{\seqNiveauGetCodeCycle}[1]{%
	\DTLfetch{param_niveaux}{Code}{#1}{CodeCycle}%
}%
% ----------------------------------------------------------------------------
% \begin{macro}{\seqThemeGetNom}%
% \IndexMacro{\seqThemeGetNom}%
% Renvoie le nom du thème spécifié dans le cycle spécifié
% #1 : code cycle, 
% #2 : code thème, 
% \begin{macrocode}%
\newcommand{\seqThemeGetNom}[2]{%
	\DTLfetch{#1_themes}{Code}{#2}{Nom}%
}%
% \end{macrocode}%
% \end{macro}%
% -------------------------------------------------------------------------------------------
% \begin{macro}{\seqThemeGetCodeCouleur}%
% \IndexMacro{\seqThemeGetCodeCouleur}%
% Renvoie le code couleur du thème spécifié dans le cycle spécifié
% #1 : code cycle, 
% #2 : code thème, 
% \begin{macrocode}%
\newcommand{\seqThemeGetCodeCouleur}[2]{%
	\DTLfetch{#1_themes}{Code}{#2}{CodeCouleur}%
}%
% \end{macrocode}%
% \end{macro}%
% -------------------------------------------------------------------------------------------
% \begin{macro}{\seqThemeGetDescription}%
% \IndexMacro{\seqThemeGetDescription}%
% Renvoie la description du thème spécifié dans le cycle spécifié
% #1 : code cycle,
% #2 : code thème,
% \begin{macrocode}%
\newcommand{\seqThemeGetDescription}[2]{%
	\DTLfetch{#1_themes}{Code}{#2}{Description}%
}%
% \end{macrocode}%
% \end{macro}%
% -------------------------------------------------------------------------------------------
% \begin{macro}{\seqSequenceGetNom}%
% \IndexMacro{\seqSequenceGetNom}%
% Renvoie le nom de la séquence courante
% \begin{macrocode}%
\newcommand{\seqSequenceGetNom}{%
	\seqDTLfetchsave{\tmp@codecycle}{param_niveaux}{Code}{\seq@code@niveau}{CodeCycle}%
	\DTLfetch{\tmp@codecycle _sequences}{Code}{\seq@code@sequence}{Nom}%
}%
% \end{macrocode}%
% \end{macro}%
% \end{macrocode}%
% \end{macro}%
% ----------------------------------------------------------------------------
% \begin{macro}{\seqSequenceGetNomOf}%
% \IndexMacro{\seqSequenceGetNomOf}%
% Variante de \seqSequenceGetNom permettant de spécifier explicitement le
% niveau et la séquence, sans modifier les variables globales courantes.
% Utile dans les tableaux récapitulatifs où l'on cite plusieurs séquences
% différentes sans changer de contexte.
% #1 : code niveau (ex. \verb|N10|)
% #2 : code séquence (ex. \verb|S03|)
% \begin{macrocode}%
\newcommand{\seqSequenceGetNomOf}[2]{%
    \seqDTLfetchsave{\tmp@codecycle}{param_niveaux}{Code}{#1}{CodeCycle}%
    \DTLfetch{\tmp@codecycle _sequences}{Code}{#2}{Nom}%
}%
% \end{macrocode}%
% \end{macro}%
% ----------------------------------------------------------------------------
% -------------------------------------------------------------------------------------------
% \begin{macro}{\seqSequenceSaveNom}%
% \IndexMacro{\seqSequenceGetNom}%
% Sauve le nom de la séquence courante dans la macro spécifiée
% #1 : macro de sauvegarde
% \begin{macrocode}%
\newcommand{\seqSequenceSaveNom}[1]{%
	\seqDTLfetchsave{\tmp@codecycle}{param_niveaux}{Code}{\seq@code@niveau}{CodeCycle}%
	\seqDTLfetchsave{#1}{\tmp@codecycle _sequences}{Code}{\seq@code@sequence}{Nom}%
}%
% \end{macrocode}%
% \end{macro}%
% -------------------------------------------------------------------------------------------
% \begin{macro}{\seqSequenceGetNumero}%
% \IndexMacro{\seqSequenceGetNumero}%
% Renvoie le numéro de la séquence spécifiée
% \begin{macrocode}%
\newcommand{\seqSequenceGetNumero}{%
	\seqDTLfetchsave{\tmp@codecycle}{param_niveaux}{Code}{\seq@code@niveau}{CodeCycle}%
	\DTLfetch{\tmp@codecycle _sequences}{Code}{\seq@code@sequence}{Numero}%
}%
% \end{macrocode}%
% \end{macro}%
% -------------------------------------------------------------------------------------------
% \begin{macro}{\seqSequenceSaveNumero}%
% \IndexMacro{\seqSequenceSaveNumero}%
% Sauve le numéro de la séquence courante dans la macro spécifiée
% #1 : macro de sauvegarde
% \begin{macrocode}%
\newcommand{\seqSequenceSaveNumero}[1]{%
	\seqDTLfetchsave{\tmp@codecycle}{param_niveaux}{Code}{\seq@code@niveau}{CodeCycle}%
	\seqDTLfetchsave{#1}{\tmp@codecycle _sequences}{Code}{\seq@code@sequence}{Numero}%
}%
% \end{macrocode}%
% \end{macro}%
% -------------------------------------------------------------------------------------------
% \begin{macro}{\seqSequenceGetTheme}%
% \IndexMacro{\seqSequenceGetTheme}%
% Renvoie le thème de la séquence courante
% \begin{macrocode}%
\newcommand{\seqSequenceGetTheme}{%
	\seqDTLfetchsave{\tmp@codecycle}{param_niveaux}{Code}{\seq@code@niveau}{CodeCycle}%
	\DTLfetch{\tmp@codecycle _sequences}{Code}{\seq@code@sequence}{Theme}%
}%
% \end{macrocode}%
% \end{macro}%
% -------------------------------------------------------------------------------------------
% \begin{macro}{\seqSequenceSaveTheme}%
% \IndexMacro{\seqSequenceSaveTheme}%
% Sauve le thème de la séquence courante dans la macro spécifiée
% #1 : macro de sauvegarde
% \begin{macrocode}%
\newcommand{\seqSequenceSaveTheme}[1]{%
	\seqDTLfetchsave{\tmp@codecycle}{param_niveaux}{Code}{\seq@code@niveau}{CodeCycle}%
	\seqDTLfetchsave{#1}{\tmp@codecycle _sequences}{Code}{\seq@code@sequence}{Theme}%
}%
% \end{macrocode}%
% \end{macro}%
% -------------------------------------------------------------------------------------------
% \begin{macro}{\seqSequenceSaveThemeOf}%
% \IndexMacro{\seqSequenceSaveThemeOf}%
% Sauve le thème de la séquence spécifiée dans la macro spécifiée
% #1 : macro de sauvegarde
% #2 : code du niveau
% #3 : code de la séquence
% \begin{macrocode}%
\newcommand{\seqSequenceSaveThemeOf}[3]{%
	\seqDTLfetchsave{\tmp@codecycle}{param_niveaux}{Code}{#2}{CodeCycle}%
	\seqDTLfetchsave{#1}{\tmp@codecycle _sequences}{Code}{#3}{Theme}%
}%
% \end{macrocode}%
% \end{macro}%
% -------------------------------------------------------------------------------------------
% \begin{macro}{\seqObjectifGetNom}%
% \IndexMacro{\seqObjectifGetNom}%
% Renvoie le nom de l'objectif spécifié dans la séquence courante
% #1 : code objectif.
% \begin{macrocode}%
\newcommand{\seqObjectifGetNom}[1]{%
	\seq@bases@filter{\seq@code@niveau}{\seq@code@sequence}%
	\DTLfetch{seq@objectifs@sequence}{CodeObjectif}{#1}{Nom}%
}%
% \end{macrocode}%
% \end{macro}%
% ----------------------------------------------------------------------------
% \begin{macro}{\seqObjectifGetFinCycle}%
% \IndexMacro{\seqObjectifGetFinCycle}%
% Renvoie si l'objectif spécifié dans la séquence courante est un objectif de fin de cycle
% #1 : code objectif.
% \begin{macrocode}%
\newcommand{\seqObjectifGetFinCycle}[1]{%
	\seq@bases@filter{\seq@code@niveau}{\seq@code@sequence}%
	\DTLfetch{seq@objectifs@sequence}{CodeObjectif}{#1}{FinCycle}%
}%
% \end{macrocode}%
% \end{macro}%
% ----------------------------------------------------------------------------
% \begin{macro}{\seqObjectifGetMaitriseTB}%
% \IndexMacro{\seqObjectifGetMaitriseTB}%
% Renvoie le critère d'atteinte du niveau Très Bon pour l'objectif spécifié dans la séquence courante
% #1 : code objectif.
% \begin{macrocode}%
\newcommand{\seqObjectifGetMaitriseTB}[1]{%
	\seq@bases@filter{\seq@code@niveau}{\seq@code@sequence}%
	\DTLfetch{seq@objectifs@sequence}{CodeObjectif}{#1}{MaitriseTB}%
}%
% \end{macrocode}%
% \end{macro}%
% ----------------------------------------------------------------------------
% \begin{macro}{\seqObjectifGetMaitriseS}%
% \IndexMacro{\seqObjectifGetMaitriseS}%
% Renvoie le critère d'atteinte du niveau Satisfaisant pour l'objectif spécifié dans la séquence courante
% #1 : code objectif, 
% \begin{macrocode}%
\newcommand{\seqObjectifGetMaitriseS}[1]{%
	\seq@bases@filter{\seq@code@niveau}{\seq@code@sequence}%
	\DTLfetch{seq@objectifs@sequence}{CodeObjectif}{#1}{MaitriseS}%
}%
% \end{macrocode}%
% \end{macro}%
% ----------------------------------------------------------------------------
% \begin{macro}{\seqObjectifGetMaitriseF}%
% \IndexMacro{\seqObjectifGetMaitriseF}%
% Renvoie le critère d'atteinte du niveau Fragile pour l'objectif spécifié dans la séquence courante
% #1 : code objectif.
% \begin{macrocode}%
\newcommand{\seqObjectifGetMaitriseF}[1]{%
	\seq@bases@filter{\seq@code@niveau}{\seq@code@sequence}%
	\DTLfetch{seq@objectifs@sequence}{CodeObjectif}{#1}{MaitriseF}%
}%
% \end{macrocode}%
% \end{macro}%
% ----------------------------------------------------------------------------
% \begin{macro}{\seqConnaissanceGetNom}%
% \IndexMacro{\seqConnaissanceGetNom}%
% Renvoie le libellé de la connaissance spécifiée dans la séquence courante
% #1 : code connaissance.
% \begin{macrocode}%
\newcommand{\seqConnaissanceGetNom}[1]{%
	\seq@bases@filter{\seq@code@niveau}{\seq@code@sequence}%
	\DTLfetch{seq@connaissances@sequence}{CodeConnaissance}{#1}{Nom}%
}%
% }%
% \end{macrocode}%
% \end{macro}%
% ----------------------------------------------------------------------------
% \begin{macro}{\seqTableauObjectifs}%
% \IndexMacro{\seqTableauObjectifs}%
% Produit le tableau récapitulatif des critères d'atteinte des objectifs
% pour la séquence courante.
% \begin{macrocode}%
\newcommand{\seqTableauObjectifs}[1][]{%
	\seq@bases@filter{\seq@code@niveau}{\seq@code@sequence}%
	\seq@rule@color{\seq@couleuri@vif}%
	\begin{center}%
		\begin{xltabular}{0.9\linewidth}{L{0.36\linewidth} *{3}{L{0.18\linewidth}}}%
			\toprule[1pt]
			\multirow{2}{*}{\centering \color{\seq@couleuri@vif}\textbf{Objectif}} &
			\multicolumn{3}{c}{\color{\seq@couleuri@vif}\textbf{Niveau de maîtrise}} \\
			\cmidrule[1pt]{2-4}%
			& \color{\seq@couleuri@vif} Très bon &
			\color{\seq@couleuri@vif} Satisfaisant &
			\color{\seq@couleuri@vif} À consolider \\
			\toprule[1pt]
			\DTLforeach*{seq@objectifs@sequence}{%
				\objCode=CodeObjectif,%
				\objNom=Nom,%
				\objMaitriseTB=MaitriseTB,%
				\objMaitriseS=MaitriseS,%
				\objMaitriseF=MaitriseF%
			}{%
				\textbf{\objCode}\newline\objNom &
				\objMaitriseTB & \objMaitriseS & \objMaitriseF \\
				\DTLiflastrow{\bottomrule[1pt]}{\midrule}%
			}%
		\end{xltabular}%
	\end{center}%
}%
% \end{macrocode}%
% \end{macro}%
% ----------------------------------------------------------------------------
% \begin{macro}{\seqTableauPrerequis}%
% \IndexMacro{\seqTableauPrerequis}%
% Produit le tableau récapitulatif des compétences à maîtriser en début 
% de séquence pour la séquence courante.
% \begin{macrocode}%
\define@cmdkey [seq]{TableauPrerequis}{niveau}{}%
\define@cmdkey [seq]{TableauPrerequis}{sequence}{}%
\newcommand{\seqTableauPrerequis}[1][]{%
	\setkeys[seq]{TableauPrerequis}{#1}%
	\ifthenelse{\equal{\cmdseq@TableauPrerequis@niveau}{N09}}{%
		% Dernière année du cycle: on reprend les objectifs des 2 années précédentes
		\seq@prerequis@annee{N07}{\cmdseq@TableauPrerequis@sequence}%
		\seq@prerequis@annee{N08}{\cmdseq@TableauPrerequis@sequence}%
	}{%
		\ifthenelse{\equal{\cmdseq@TableauPrerequis@niveau}{N10}}{%
			% Première année du cycle: on reprend les attendus de fin du cycle précédent
			\seq@prerequis@attenduscycleprecedent{\cmdseq@TableauPrerequis@niveau}{\cmdseq@TableauPrerequis@sequence}%
		}{%
			\ifthenelse{\equal{\cmdseq@TableauPrerequis@niveau}{N11}}{%
				% Deuxième année du cycle: on reprend les objectif de l'année précédente
				\seq@prerequis@annee{N10}{\cmdseq@TableauPrerequis@sequence}%
			}{%
				\ifthenelse{\equal{\cmdseq@TableauPrerequis@niveau}{N12}}{%
					% Dernière année du cycle: on reprend les objectifs des 2 années précédentes
					\seq@prerequis@annee{N10}{\cmdseq@TableauPrerequis@sequence}%
					\seq@prerequis@annee{N11}{\cmdseq@TableauPrerequis@sequence}%
				}{%
					% code niveau inconnu : générer une erreur
				}%
			}%
		}%
	}%
}%
% \end{macrocode}%
% \end{macro}%
% ----------------------------------------------------------------------------
% \begin{macro}{\seqMaitriseGetNom}%
% \IndexMacro{\seqMaitriseGetNom}%
% Renvoie le nom d'un niveau de maîtrise spécifié.
% #1 : code du niveau de maîtrise
% \begin{macrocode}%
\newcommand{\seqMaitriseGetNom}[1]{%
	\ifthenelse{\equal{#1}{I}}{Insuffisant}{%
		\ifthenelse{\equal{#1}{F}}{\`A consolider}{%
			\ifthenelse{\equal{#1}{A}}{Satisfaisant}{%
				\ifthenelse{\equal{#1}{E}}{Très bon}{%
					\ifthenelse{\equal{#1}{NE}}{Non évalué}{%
						#1 % inconnu: on renvoie la même chose
					}%
				}%
			}%
		}%
	}%
}%
% \end{macrocode}%
% \end{macro}%
% ----------------------------------------------------------------------------
% \begin{macro}{\seqSuiviAnnuelParEleve}%
% \IndexMacro{\seqSuiviAnnuelParEleve}%
% Produit le tableau annuel de suivi des compétences par élève 
% pour une liste d'élèves.
% \begin{macrocode}%
\define@cmdkey[seq]{suiviAnnuelParEleve}{niveau}{}%
\newcommand{\seqSuiviAnnuelParEleve}[1][]{%
	\setkeys[seq]{suiviAnnuelParEleve}{#1}%
	% On reconstruit les bases objectifs par séquence
	\seq@basesobjectifs@filter@niveau{\cmdseq@suiviAnnuelParEleve@niveau}%
	% On traite les données
	\DTLforeach*{eleves}{\tmp@numEleve=Numero,\tmp@nomEleve=Nom,\tmp@prenomEleve=Prenom}{%
		\def\tblrbody{}%
		\edef\tblrbody{\tblrbody \SetCell[c=2]{c} \tmp@nomEleve\;\tmp@prenomEleve & \\}%
		\DTLforeach*{periodes}{\tmp@periode=Periode}{%
			\edef\tblrbody{\tblrbody \SetRow{bg=blue!30!white} Bilan\;\tmp@periode \\}%
			\DTLforeach*{sequencesPeriode\tmp@periode}{\tmp@codeSequence=Code,\tmp@nomSequence=Nom}{%
				\edef\tblrbody{\tblrbody \tmp@codeSequence\;-\;\tmp@nomSequence & }%
				\seqDTLfetchsave{\tmp@niveauMaitrise}{bilanPeriode\tmp@periode}{Numero}{\tmp@numEleve}{\tmp@codeSequence}%
				\seq@tblrcell@decodemaitrise{\tmp@niveauMaitrise} \edef\tblrbody{\tblrbody \\}%
			}%
		}%
		\seqDTLfetchsave{\tmp@annee}{param_niveaux}{Code}{\cmdseq@suiviAnnuelParEleve@niveau}{AnneeDansCycle}%
		\DTLforeach*{C04_sequences}{\tmp@codeSequence=Code,\tmp@nomSequence=Nom,\tmp@codeTheme=Theme}{%
			\edef\tblrbody{\tblrbody \SetRow{bg=blue!30!white} \tmp@codeSequence\;-\;\tmp@nomSequence \\}%
			\seq@seqcode@tobasename{\tmp@nombase}{\tmp@codeSequence}%
			\DTLforeach*{\tmp@nombase}{%
					\tmp@codeObjectif=CodeObjectif,\tmp@nomObjectif=Nom}{%
				\edef\tblrbody{\tblrbody \tmp@nomObjectif &}%
				\seqDTLfetchsave{\tmp@niveauMaitrise}{Suivi\tmp@codeSequence}{Numero}{\tmp@numEleve}{\tmp@codeObjectif}%
				\seq@tblrcell@decodemaitrise{\tmp@niveauMaitrise} \edef\tblrbody{\tblrbody \\}%
			}%
		}%
		\begin{longtblr}[expand=\tblrbody]{%
			width=0.95\linewidth,colspec={X[4]X[c]},hlines,
			rows={valign=m},
			rowhead=1
			}%
			\tblrbody
		\end{longtblr}%
		\cleardoublepage
	}%
}%
% \end{macrocode}%
% \end{macro}%
% ----------------------------------------------------------------------------
% \begin{macro}{\seqSuiviAnnuelParSequence}%
% \IndexMacro{\seqSuiviAnnuelParSequence}%
% Produit le tableau de suivi des compétences d'une séquence.
% pour une liste d'élèves.
% \begin{macrocode}%
\define@cmdkey[seq]{suiviAnnuelParSequence}{niveau}{}%
\define@cmdkey[seq]{suiviAnnuelParSequence}{sequence}{}%
\newcommand{\seqSuiviAnnuelParSequence}[1][]{%
	%% ---- Options
	\setkeys[seq]{suiviAnnuelParSequence}{#1}%
	% On reconstruit les bases objectifs par séquence
	\seq@basesobjectifs@filter@niveau{\cmdseq@suiviAnnuelParSequence@niveau}%
	% On traite les données
	\def\tblrbody{}%
	\DTLforeach*{eleves}{\tmp@nomEleve=Nom,\tmp@prenomEleve=Prenom}{%
		\edef\tblrbody{\tblrbody & \rotatebox{90}{\tmp@nomEleve\;\tmp@prenomEleve}}%
	}%
	\edef\tblrbody{\tblrbody \\}%
	\seq@seqcode@tobasename{\tmp@nombase}{\cmdseq@suiviAnnuelParSequence@sequence}%
	\DTLforeach*{\tmp@nombase}{\tmp@codeObj=CodeObjectif,\tmp@nomObj=Nom}{%
		\edef\tblrbody{\tblrbody Objectif \tmp@codeObj \newline \tmp@nomObj &}%
		\DTLforeach*{eleves}{\tmp@numEleve=Numero,\tmp@nomEleve=Nom,\tmp@prenomEleve=Prenom}{%
			\seqDTLfetchsave{\tmp@niveauMaitrise}{Suivi\cmdseq@suiviAnnuelParSequence@sequence}{Numero}{\tmp@numEleve}{\tmp@codeObj}%
			\seq@tblrcell@coloremaitrise{\tmp@niveauMaitrise} \edef\tblrbody{\tblrbody \tmp@niveauMaitrise}%
			\DTLiflastrow{}{\edef\tblrbody{\tblrbody &}}%
		}%
		\edef\tblrbody{\tblrbody \\}%
	}%
	\begin{center}%
		\begin{longtblr}[expand=\tblrbody]{%
				width=0.9\linewidth,colspec={X[8]*{25}{X[c]}},hlines,vlines,
				rows={valign=m},columns={font={\bfseries}},column{1}={font={\large}},row{1}={font={\normalsize}}}%
			\tblrbody
		\end{longtblr}%
	\end{center}%
}%
% \end{macrocode}%
% \end{macro}%
% ----------------------------------------------------------------------------
% \begin{macro}{\seqSuiviPeriodiqueParSequence}%
% \IndexMacro{\seqSuiviPeriodiqueParSequence}%
% Produit le tableau de suivi des compétences des séquences d'une période
% pour une liste d'élèves.
% \begin{macrocode}%
\define@cmdkey[seq]{suiviPeriodiqueParSequence}{periode}{}%
\define@cmdkey[seq]{suiviPeriodiqueParSequence}{niveau}{}%
\newcommand{\seqSuiviPeriodiqueParSequence}[1][]{%
	%% ---- Options
	\setkeys[seq]{suiviPeriodiqueParSequence}{#1}%
	% On reconstruit les bases objectifs par séquence
	\seq@basesobjectifs@filter@niveau{\cmdseq@suiviPeriodiqueParSequence@niveau}%
	% On traite les données
	\def\tblrbody{}%
	\DTLforeach*{eleves}{\tmp@nomEleve=Nom,\tmp@prenomEleve=Prenom}{%
		\edef\tblrbody{\tblrbody & \rotatebox{90}{\tmp@nomEleve\;\tmp@prenomEleve}}%
	}%
	\edef\tblrbody{\tblrbody \\}%
	\DTLforeach*{sequencesPeriode\cmdseq@suiviPeriodiqueParSequence@periode}{\tmp@codesequence=Code,\tmp@nomsequence=Nom}{%
		\edef\tblrbody{\tblrbody \tmp@codesequence\;-\;\tmp@nomsequence & }%
		\DTLforeach*{eleves}{\tmp@numEleve=Numero}{%
			\seqDTLfetchsave{\tmp@niveauMaitrise}{bilanPeriode\cmdseq@suiviPeriodiqueParSequence@periode}{Numero}{\tmp@numEleve}{\tmp@codesequence}%
			\seq@tblrcell@coloremaitrise{\tmp@niveauMaitrise} \edef\tblrbody{\tblrbody \tmp@niveauMaitrise}%
			\DTLiflastrow{}{\edef\tblrbody{\tblrbody &}}%
		}%
		\edef\tblrbody{\tblrbody \\}%
	}%
	\begin{center} \begin{longtblr}[expand=\tblrbody]{%
			width=0.9\linewidth,colspec={X[8]*{25}{X[c]}},hlines,vlines,
			rows={valign=m},columns={font={\bfseries}},column{1}={font={\large}},row{1}={font={\normalsize}}}%
			\tblrbody
		\end{longtblr} \end{center}%
	\bigskip\par
	\centerline{\Large{\phantom{test}\hfill Dates des séquences \hfill}}%
	\medskip\par
	\DTLforeach*{sequencesPeriode\cmdseq@suiviPeriodiqueParSequence@periode}{\tmp@codeSequence=Code,\tmp@nomSequence=Nom,\tmp@dateDebut=Debut,\tmp@dateFin=Fin}{%
		\large{\phantom{test}\qquad\tmp@codeSequence\;-\;\tmp@nomSequence: du \tmp@dateDebut{} au \tmp@dateFin}\smallskip\par
	}%
}%
% \end{macrocode}%
% \end{macro}%
% ----------------------------------------------------------------------------
% \begin{macro}{\seqPlanReperesNiveaux}%
% \IndexMacro{\seqPlanReperesNiveaux}%
% Produit le tableau des repères temporels (date début et fin) et 
% la répartition des séances de la séquence.
% \begin{macrocode}%
\define@cmdkey[seq]{blocReperesNiveaux}{objCours}[01]{}%
\define@cmdkey[seq]{blocReperesNiveaux}{objExosFirst}[02]{}%
\define@cmdkey[seq]{blocReperesNiveaux}{objExosLast}[09]{}%
\define@cmdkey[seq]{blocReperesNiveaux}{dateDeb}[\ldots/\ldots]{}%
\define@cmdkey[seq]{blocReperesNiveaux}{dateFin}[\ldots/\ldots]{}%
\define@cmdkey[seq]{blocReperesNiveaux}{seancesTotal}[\ldots]{}%
\define@cmdkey[seq]{blocReperesNiveaux}{seancesRevisions}[\ldots]{}%
\define@cmdkey[seq]{blocReperesNiveaux}{seancesCours}[\ldots]{}%
\define@cmdkey[seq]{blocReperesNiveaux}{seancesObj02}[\ldots]{}%
\define@cmdkey[seq]{blocReperesNiveaux}{seancesObj03}[\ldots]{}%
\define@cmdkey[seq]{blocReperesNiveaux}{seancesObj04}[\ldots]{}%
\define@cmdkey[seq]{blocReperesNiveaux}{seancesObj05}[\ldots]{}%
\define@cmdkey[seq]{blocReperesNiveaux}{seancesObj06}[\ldots]{}%
\define@cmdkey[seq]{blocReperesNiveaux}{seancesObj12}[\ldots]{}%
\define@cmdkey[seq]{blocReperesNiveaux}{seancesObj13}[\ldots]{}%
\define@cmdkey[seq]{blocReperesNiveaux}{seancesObj14}[\ldots]{}%
\define@cmdkey[seq]{blocReperesNiveaux}{seancesObj15}[\ldots]{}%
\define@cmdkey[seq]{blocReperesNiveaux}{seancesObj16}[\ldots]{}%
\define@cmdkey[seq]{blocReperesNiveaux}{seancesObj22}[\ldots]{}%
\define@cmdkey[seq]{blocReperesNiveaux}{seancesObj23}[\ldots]{}%
\define@cmdkey[seq]{blocReperesNiveaux}{seancesObj24}[\ldots]{}%
\define@cmdkey[seq]{blocReperesNiveaux}{seancesObj25}[\ldots]{}%
\define@cmdkey[seq]{blocReperesNiveaux}{seancesObj26}[\ldots]{}%
\newcommand{\seqPlanReperesNiveaux}[1][]{%
	\setkeys[seq]{blocReperesNiveaux}{objCours,objExosFirst,objExosLast,dateDeb,dateFin,seancesTotal,seancesRevisions,seancesCours,#1}%
	\seq@bases@filter{\seq@code@niveau}{\seq@code@sequence}%
	\begin{multicols}{2}%
		\begin{blocReperesNiveaux}{Repères temporels}%
			\item Dates: du \cmdseq@blocReperesNiveaux@dateDeb~au~\cmdseq@blocReperesNiveaux@dateFin
			\item Travail en classe: \cmdseq@blocReperesNiveaux@seancesTotal~séances
			\begin{seqColItem}[nbCols=1,label=\ding{51}]
				\item Révisions: \cmdseq@blocReperesNiveaux@seancesRevisions~séances
				\item Cours: \cmdseq@blocReperesNiveaux@seancesCours~séances
				\DTLforeach[\DTLisiclosedbetween{\codeObjectif}{\cmdseq@blocReperesNiveaux@objExosFirst}{\cmdseq@blocReperesNiveaux@objExosLast}]
				{seq@objectifs@sequence}{\codeObjectif=CodeObjectif}{%
					\item Objectif \codeObjectif: \csname cmdseq@blocReperesNiveaux@seancesObj\codeObjectif\endcsname~séances
				}%
			\end{seqColItem}%
		\end{blocReperesNiveaux}%
		\columnbreak
		\begin{blocReperesNiveaux}{Niveaux atteints}%
			\DTLforeach[\DTLisiclosedbetween{\codeObjectif}{\cmdseq@blocReperesNiveaux@objCours}{\cmdseq@blocReperesNiveaux@objExosLast}]
			{seq@objectifs@sequence}{\codeObjectif=CodeObjectif}{%
				\item Objectif \codeObjectif: \ldots
			}%
			\item \begin{minipage}[c][3em][c]{\linewidth} \textbf{Note équivalente: \ldots/20} \end{minipage}%
		\end{blocReperesNiveaux}%
	\end{multicols}%
}%
% \end{macrocode}%
% \end{macro}%
% ============================================================================
% Macros d'itération sur thèmes et séquences (Axe 3 – documents annuels)
% ============================================================================
% Ces macros encapsulent datatool pour les documents qui ont besoin d'itérer
% sur l'ensemble des thèmes et des séquences du cycle courant (cours, exercices).
% Elles lisent le cycle depuis \seq@code@niveau via la base param_niveaux.
%
% \begin{macro}{\seqForeachTheme}%
% \IndexMacro{\seqForeachTheme}%
% Itère sur tous les thèmes du cycle correspondant au niveau courant.
% #1 : macro qui recevra le code thème
% #2 : macro qui recevra le nom du thème
% #3 : macro qui recevra la description du thème
% #4 : macro qui recevra le code couleur du thème
% #5 : corps à exécuter pour chaque thème
% \begin{macrocode}%
\newcommand{\seqForeachTheme}[5]{%
    \seqDTLfetchsave{\seq@iter@cycle}{param_niveaux}{Code}{\seq@code@niveau}{CodeCycle}%
    \DTLforeach*{\seq@iter@cycle _themes}{%
        #1=Code,%
        #2=Nom,%
        #3=Description,%
        #4=CodeCouleur%
    }{#5}%
}%
% \end{macrocode}%
% \end{macro}%
% ----------------------------------------------------------------------------
% \begin{macro}{\seqForeachSequence}%
% \IndexMacro{\seqForeachSequence}%
% Itère sur toutes les séquences du cycle correspondant au niveau courant.
% #1 : macro qui recevra le code séquence
% #2 : macro qui recevra le nom de la séquence
% #3 : macro qui recevra le code thème de la séquence
% #4 : corps à exécuter pour chaque séquence
% \begin{macrocode}%
\newcommand{\seqForeachSequence}[4]{%
    \seqDTLfetchsave{\seq@iter@cycle}{param_niveaux}{Code}{\seq@code@niveau}{CodeCycle}%
    \DTLforeach*{\seq@iter@cycle _sequences}{%
        #1=Code,%
        #2=Nom,%
        #3=Theme%
    }{#4}%
}%
% \end{macrocode}%
% \end{macro}%
% ----------------------------------------------------------------------------
% Variable interne temporaire pour les macros d'itération
\newcommand{\seq@iter@cycle}{}%
% ----------------------------------------------------------------------------

% -------------------------------------------------------------------------------------------
% \begin{macro}{\seqSetColors}%
% \IndexMacro{\seqSetColors}%
% Règle les couleurs de la séquence courante en appelant \seq@setColors
% \begin{macrocode}%
\newcommand{\seqSetColors}{%
	\ifthenelse{\equal{\seq@code@niveau}{N00}}{\seq@setColors{gris}}{%
		\seqDTLfetchsave{\tmp@codecycle}{param_niveaux}{Code}{\seq@code@niveau}{CodeCycle}%
		\seqDTLfetchsave{\tmp@codeTheme}{\tmp@codecycle _sequences}{Code}{\seq@code@sequence}{Theme}%
		\seqDTLfetchsave{\tmp@codeCouleur}{\tmp@codecycle _themes}{Code}{\tmp@codeTheme}{CodeCouleur}%
		\seq@setColors{\tmp@codeCouleur}%
	}%
}%
% \end{macrocode}%
% \end{macro}%
% ----------------------------------------------------------------------------

% ----------------------------------------------------------------------------
% \begin{macro}{\seqSetColorsTheme}%
% \IndexMacro{\seqSetColorsTheme}%
% Wrapper public pour \seq@setColors
% #1 : code de couleur à appliquer
% \begin{macrocode}%
\newcommand{\seqSetColorsTheme}[1]{%
	\seq@setColors{#1}%
}%
% \end{macrocode}%
% \end{macro}%
% ----------------------------------------------------------------------------

% ----------------------------------------------------------------------------
% \begin{macro}{\seqFiltreDonnees}%
% \IndexMacro{\seqSetColorsTheme}%
% Wrapper public pour \seq@bases@filter avec la séquence courante
% \begin{macrocode}%
\newcommand{\seqFiltreDonnees}{%
	\seq@bases@filter{\seq@code@niveau}{\seq@code@sequence}%
}%
% \end{macrocode}%
% \end{macro}%
% ----------------------------------------------------------------------------
%
% ============================================================================
% Commandes internes
% ============================================================================
% ----------------------------------------------------------------------------
% Vérifie que les bases sont filtrées pour le niveau et la séquence spécifiés
% #1 : code de niveau
% #2 : code de séquence
\newcommand{\seq@bases@filter}[2]{%
	\seqDTLfetchsave{\tmp@codecycle}{param_niveaux}{Code}{#1}{CodeCycle}%
	% On vérifie les bases filtrées par cycle
	\ifthenelse{\equal{\tmp@codecycle}{C03}}{%
		\ifthenelse{\equal{\seq@var@cycle}{C03}}{%cache cycle 3 valide : rien à faire
		}{%mise à jour des bases filtrées pour le cycle 3
			\gdef\seq@var@cycle{C03}%
			\seq@bases@filter@cycle{C03}{N07}{N08}{N09}%
		}%
	}{%
		\ifthenelse{\equal{\tmp@codecycle}{C04}}{%
			\ifthenelse{\equal{\seq@var@cycle}{C04}}{%cache cycle 4 valide : rien à faire
			}{%mise à jour des bases filtrées pour le cycle 4
				\gdef\seq@var@cycle{C04}%
				\seq@bases@filter@cycle{C04}{N10}{N11}{N12}%
			}%
		}{%cycle inconnu: générer une erreur
		}%
	}%
	% On vérifie les bases filtrées par niveau et séquence 
	\ifthenelse{\equal{#1}{\seq@var@niveau}\and\equal{#2}{\seq@var@sequence}}{%cache niveau/séquence valide : rien à faire
	}{%mise à jour des bases filtrées pour #2 de #1
		\xdef\seq@var@niveau{#1}%
		\xdef\seq@var@sequence{#2}%
		\seqDTLfetchsave{\tmp@annee}{param_niveaux}{Code}{#1}{AnneeDansCycle}%
		\seq@bases@filter@sequence{seq@connaissances@\tmp@annee}{seq@objectifs@\tmp@annee}{#2}%
	}%
}%
% ----------------------------------------------------------------------------
% Reconstruit les bases connaissances et objectifs filtrées 
% par niveau pour le cycle spécifié
% #1 : code cycle
% #2 : code niveau un
% #3 : code niveau deux
% #4 : code niveau trois
\newcommand{\seq@bases@filter@cycle}[4]{%
	\DTLcleardb{seq@connaissances@anneeun}%
	\DTLcleardb{seq@connaissances@anneedeux}%
	\DTLcleardb{seq@connaissances@anneetrois}%
	\DTLforeach*{#1_connaissances}{%
		\CodeNiveau=CodeNiveau,%
		\CodeSequence=CodeSequence,%
		\CodeConnaissance=CodeConnaissance,%
		\Nom=Nom%
	}{%
		\dtlexpandnewvalue
		\DTLifstringeq{\CodeNiveau}{#2}{%
			\DTLnewrow{seq@connaissances@anneeun}%
			\DTLnewdbentry{seq@connaissances@anneeun}{CodeSequence}{\CodeSequence}%
			\DTLnewdbentry{seq@connaissances@anneeun}{CodeConnaissance}{\CodeConnaissance}%
			\DTLnewdbentry{seq@connaissances@anneeun}{Nom}{\Nom}%
		}{%
			\DTLifstringeq{\CodeNiveau}{#3}{%
				\DTLnewrow{seq@connaissances@anneedeux}%
				\DTLnewdbentry{seq@connaissances@anneedeux}{CodeSequence}{\CodeSequence}%
				\DTLnewdbentry{seq@connaissances@anneedeux}{CodeConnaissance}{\CodeConnaissance}%
				\DTLnewdbentry{seq@connaissances@anneedeux}{Nom}{\Nom}%
			}{%
				\DTLifstringeq{\CodeNiveau}{#4}{%
					\DTLnewrow{seq@connaissances@anneetrois}%
					\DTLnewdbentry{seq@connaissances@anneetrois}{CodeSequence}{\CodeSequence}%
					\DTLnewdbentry{seq@connaissances@anneetrois}{CodeConnaissance}{\CodeConnaissance}%
					\DTLnewdbentry{seq@connaissances@anneetrois}{Nom}{\Nom}%
				}{%
				}%
			}%
		}%
	}%
	\DTLcleardb{seq@objectifs@anneeun}%
	\DTLcleardb{seq@objectifs@anneedeux}%
	\DTLcleardb{seq@objectifs@anneetrois}%
	\DTLforeach*{#1_objectifs}{%
		\CodeNiveau=CodeNiveau,%
		\CodeSequence=CodeSequence,%
		\CodeObjectif=CodeObjectif,%
		\Nom=Nom,%
		\MaitriseTB=MaitriseTB,%
		\MaitriseS=MaitriseS,%
		\MaitriseF=MaitriseF%
	}{%
		\dtlexpandnewvalue
		\DTLifstringeq{\CodeNiveau}{#2}{%
			\DTLnewrow{seq@objectifs@anneeun}%
			\DTLnewdbentry{seq@objectifs@anneeun}{CodeSequence}{\CodeSequence}%
			\DTLnewdbentry{seq@objectifs@anneeun}{CodeObjectif}{\CodeObjectif}%
			\DTLnewdbentry{seq@objectifs@anneeun}{Nom}{\Nom}%
			\DTLnewdbentry{seq@objectifs@anneeun}{MaitriseTB}{\MaitriseTB}%
			\DTLnewdbentry{seq@objectifs@anneeun}{MaitriseS}{\MaitriseS}%
			\DTLnewdbentry{seq@objectifs@anneeun}{MaitriseF}{\MaitriseF}%
		}{%
			\DTLifstringeq{\CodeNiveau}{#3}{%
				\DTLnewrow{seq@objectifs@anneedeux}%
				\DTLnewdbentry{seq@objectifs@anneedeux}{CodeSequence}{\CodeSequence}%
				\DTLnewdbentry{seq@objectifs@anneedeux}{CodeObjectif}{\CodeObjectif}%
				\DTLnewdbentry{seq@objectifs@anneedeux}{Nom}{\Nom}%
				\DTLnewdbentry{seq@objectifs@anneedeux}{MaitriseTB}{\MaitriseTB}%
				\DTLnewdbentry{seq@objectifs@anneedeux}{MaitriseS}{\MaitriseS}%
				\DTLnewdbentry{seq@objectifs@anneedeux}{MaitriseF}{\MaitriseF}%
			}{%
				\DTLifstringeq{\CodeNiveau}{#4}{%
					\DTLnewrow{seq@objectifs@anneetrois}%
					\DTLnewdbentry{seq@objectifs@anneetrois}{CodeSequence}{\CodeSequence}%
					\DTLnewdbentry{seq@objectifs@anneetrois}{CodeObjectif}{\CodeObjectif}%
					\DTLnewdbentry{seq@objectifs@anneetrois}{Nom}{\Nom}%
					\DTLnewdbentry{seq@objectifs@anneetrois}{MaitriseTB}{\MaitriseTB}%
					\DTLnewdbentry{seq@objectifs@anneetrois}{MaitriseS}{\MaitriseS}%
					\DTLnewdbentry{seq@objectifs@anneetrois}{MaitriseF}{\MaitriseF}%
				}{%
				}%
			}%
		}%
	}%
}%
% ----------------------------------------------------------------------------
% Reconstruit les bases objectifs et connaissances filtrées 
% pour la séquence spécifiée
% #1 : BBD connaissances du niveau
% #2 : BBD objectifs du niveau
% #3 : code séquence
\newcommand{\seq@bases@filter@sequence}[3]{%
	% connaissances
	\DTLcleardb{seq@connaissances@sequence}%
	\DTLforeach*[\DTLiseq{\CodeSequence}{#3}]{#1}{%
		\CodeSequence=CodeSequence,%
		\CodeConnaissance=CodeConnaissance,%
		\Nom=Nom%
	}{%
		\dtlexpandnewvalue
		\DTLnewrow{seq@connaissances@sequence}%
		\DTLnewdbentry{seq@connaissances@sequence}{CodeConnaissance}{\CodeConnaissance}%
		\DTLnewdbentry{seq@connaissances@sequence}{Nom}{\Nom}%
	}%
	% objectifs
	\DTLcleardb{seq@objectifs@sequence}%
	\DTLforeach*[\DTLiseq{\CodeSequence}{#3}]{#2}{%
		\CodeSequence=CodeSequence,%
		\CodeObjectif=CodeObjectif,%
		\Nom=Nom,%
		\MaitriseTB=MaitriseTB,%
		\MaitriseS=MaitriseS,%
		\MaitriseF=MaitriseF%
	}{%
		\dtlexpandnewvalue
		\DTLnewrow{seq@objectifs@sequence}%
		\DTLnewdbentry{seq@objectifs@sequence}{CodeObjectif}{\CodeObjectif}%
		\DTLnewdbentry{seq@objectifs@sequence}{Nom}{\Nom}%
		\DTLnewdbentry{seq@objectifs@sequence}{MaitriseTB}{\MaitriseTB}%
		\DTLnewdbentry{seq@objectifs@sequence}{MaitriseS}{\MaitriseS}%
		\DTLnewdbentry{seq@objectifs@sequence}{MaitriseF}{\MaitriseF}%
	}%
}%
% ----------------------------------------------------------------------------
% ----------------------------------------------------------------------------
% Reconstruit les bases objectifs filtrées 
% par séquence pour le niveau spécifié
% #1 : code niveau
\newcommand{\seq@basesobjectifs@filter@niveau}[1]{%
	\ifthenelse{\equal{\seq@var@niveau}{#1}}{%cache niveau valide : rien à faire
	}{%mise à jour des bases filtrées pour le niveau
		\xdef\seq@var@niveau{#1}%
		% Vider toutes les bases par séquence en boucle
		% (couvre S01 à S\seq@maxsequences, quel que soit le nombre réel de séquences)
		\setcounter{seq@tmp@i}{1}%
		\@whilenum\value{seq@tmp@i}<\numexpr\seq@maxsequences+1\relax\do{%
			\DTLcleardb{seq@objectifs@s\arabic{seq@tmp@i}}%
			\stepcounter{seq@tmp@i}%
		}%
		\seqDTLfetchsave{\tmp@annee}{param_niveaux}{Code}{#1}{AnneeDansCycle}%
		\DTLforeach*{seq@objectifs@\tmp@annee}{%
			\CodeSequence=CodeSequence,%
			\CodeObjectif=CodeObjectif,%
			\Nom=Nom,%
			\MaitriseTB=MaitriseTB,%
			\MaitriseS=MaitriseS,%
			\MaitriseF=MaitriseF%
		}{%
			% Calcul direct du nom de base : S01 → seq@objectifs@s1, S12 → seq@objectifs@s12
			\seq@seqcode@tonumber{\tmp@seqnum}{\CodeSequence}%
			\edef\tmp@base{seq@objectifs@s\tmp@seqnum}%
			\dtlexpandnewvalue
			\DTLnewrow{\tmp@base}%
			\DTLnewdbentry{\tmp@base}{CodeObjectif}{\CodeObjectif}%
			\DTLnewdbentry{\tmp@base}{Nom}{\Nom}%
			\DTLnewdbentry{\tmp@base}{MaitriseTB}{\MaitriseTB}%
			\DTLnewdbentry{\tmp@base}{MaitriseS}{\MaitriseS}%
			\DTLnewdbentry{\tmp@base}{MaitriseF}{\MaitriseF}%
		}%
	}%
}%
% ----------------------------------------------------------------------------
% \begin{macro}{\seq@seqcode@tobasename}%
% Remplace l'ancienne \seq@sequence@save@nombaseobjectifs.
% Calcule le nom de la base objectifs d'une séquence à partir de son code,
% sans cascade de comparaisons : S01 → seq@objectifs@s1, S16 → seq@objectifs@s16.
% Fonctionne pour tout N entre 1 et \seq@maxsequences.
% #1 : macro de sauvegarde
% #2 : code séquence (ex. S01, S12)
\newcommand{\seq@seqcode@tobasename}[2]{%
	\seq@seqcode@tonumber{\seq@tmp@n}{#2}%
	\edef#1{seq@objectifs@s\seq@tmp@n}%
}%
% ----------------------------------------------------------------------------
% Récapitule les attendus de fin de cycle du cycle précédent 
% pour la séquence spécifiée
% #1 : code niveau
% #2 : code séquence
\newcommand{\seq@prerequis@attenduscycleprecedent}[2]{%
	\ifthenelse{\equal{#1}{N10}}{%
		\begin{seq@prerequis@N1}[Les connaissances et savoir-faire suivants sont l'ensemble des attendus de fin de cycle 3]
			\DTLforeach*[\DTLiseq{\dtl@codeSequence}{#2}]{C04_attendus_fin_C03}{\dtl@codeSequence=CodeSequence,\attendu=Attendu}{\attendu}%
		\end{seq@prerequis@N1}%
	}{% code cycle inconnu : générer une erreur
	}%
}%
% ----------------------------------------------------------------------------
% Récapitule les objectifs de l'année pour la séquence spécifiée
% #1 : code niveau
% #2 : code séquence
\newcommand{\seq@prerequis@annee}[2]{%
	\seq@bases@filter{#1}{#2}%
	\begin{seq@prerequis@N1}[Les connaissances et savoir-faire suivants sont l'ensemble des attendus de fin de \seqNiveauGetNomCourt{#1}]
		\edef\seq@tmp@rowcount{\DTLrowcount{seq@connaissances@sequence}}%
		\DTLifeq{\seq@tmp@rowcount}{0}{% Pas de nouvelle connaissance
			\item Aucune nouvelle connaissance pour cette séquence.
		}{%
			\item Connaissances de la séquence:
				\begin{seq@prerequis@N2}%
					\DTLforeach*{seq@connaissances@sequence}{\nom=Nom}{\item \nom}%
				\end{seq@prerequis@N2}%
		}%
		\edef\seq@tmp@rowcount{\DTLrowcount{seq@objectifs@sequence}}%
		\DTLifeq{\seq@tmp@rowcount}{0}{% Pas de nouvel objectif
			\item Aucun nouvel objectif pour cette séquence.
		}{%
			\item Objectifs de la séquence:
				\begin{seq@prerequis@N2}%
					\DTLforeach*{seq@objectifs@sequence}{\nom=Nom}{\item \nom}%
				\end{seq@prerequis@N2}%
		}%
	\end{seq@prerequis@N1}%
}%
% ---------------------------------------------------------------------------------
% Formatte une cellule dans un tblr pour le niveau de maîtrise spécifié
% #1 : code niveau de maîtrise
\newcommand{\seq@tblrcell@decodemaitrise}[1]{%
	\DTLifeq{#1}{I}{\edef\tblrbody{\tblrbody \SetCell{bg=red,fg=white} Insuffisant}}{%
		\DTLifeq{#1}{F}{\edef\tblrbody{\tblrbody \SetCell{bg=yellow,fg=black} \`A consolider}}{%
			\DTLifeq{#1}{A}{\edef\tblrbody{\tblrbody \SetCell{bg=green!60!white,fg=black} Satisfaisant}}{%
				\DTLifeq{#1}{E}{\edef\tblrbody{\tblrbody \SetCell{bg=green!60!black,fg=white} Très bon}}{%
					\DTLifeq{#1}{NE}{\edef\tblrbody{\tblrbody Non évalué}}{%
						\DTLifnumerical{#1}{\edef\tblrbody{\tblrbody #1 }}{%
							\DTLifeq{#1}{-}{\edef\tblrbody{\tblrbody -}}{\edef\tblrbody{\tblrbody INCONNU:#1}}}}}}}}%
}%
% \end{macro}%
% ---------------------------------------------------------------------------------
% Colore une cellule vide dans un tblr pour le niveau de maîtrise spécifié
% #1 : code niveau de maîtrise
% \begin{macro}{seqColoreMaitriseTblr}%
\newcommand{\seq@tblrcell@coloremaitrise}[1]{%
	\DTLifeq{#1}{I}{\edef\tblrbody{\tblrbody \SetCell{bg=red,fg=white}}}{%
		\DTLifeq{#1}{F}{\edef\tblrbody{\tblrbody \SetCell{bg=yellow,fg=black}}}{%
			\DTLifeq{#1}{A}{\edef\tblrbody{\tblrbody \SetCell{bg=green!60!white,fg=black}}}{%
				\DTLifeq{#1}{E}{\edef\tblrbody{\tblrbody \SetCell{bg=green!60!black,fg=white}}}{%
					\DTLifeq{#1}{NE}{}{%
						\DTLifeq{#1}{-}{}{\edef\tblrbody{\tblrbody \SetCell{bg=gray!50!white,fg=white}}}}}}}}%
}%
% \end{macro}%
% ---------------------------------------------------------------------------------
%
% ============================================================================
% Environnements internes
% ============================================================================
% Met en forme le premier niveau du tableau de prérequis
% #1 : Commentaire explicatif du contenu du tableau
\newenvironment{seq@prerequis@N1}[1][Les connaissances et savoir-faire suivants sont
	nécessaires pour réussir la séquence.]{%
	\par\vspace{20pt}%
	\begin{boitePaleNoBreak}{Prérequis}%
		#1
		\tcblower
		\begin{minipage}{\linewidth}%
			\begin{multicols}{2} \begin{dinglist}{224}%
					}{%
				\end{dinglist} \end{multicols}%
		\end{minipage}%
	\end{boitePaleNoBreak}%
	\par\vspace{10pt}%
}%
% \end{environment}%
% ----------------------------------------------------------------------------
% Met en forme le second niveau du tableau de prérequis
\newenvironment{seq@prerequis@N2}{%
	\begin{dinglist}{118}%
		}{%
	\end{dinglist}%
}%
% ----------------------------------------------------------------------------
% Fin du module data
%</package>
